\documentclass[runningheads]{llncs}

\usepackage[T1]{fontenc}
\usepackage{amsmath,amssymb}
\usepackage{graphicx}
\usepackage{booktabs}
\usepackage{multirow}
\usepackage{subcaption}
\usepackage{array}
\usepackage{tabularx}
\usepackage{xcolor}
\usepackage{url}
\usepackage{enumitem}
\usepackage{algorithm}
\usepackage{algpseudocode}

\newcommand{\concat}{\mathbin{\Vert}}
\newcommand{\sigmoid}{\sigma}
\newcommand{\placeholderfigure}[2]{%
\fbox{\begin{minipage}[c][#2][c]{0.94\textwidth}\centering\small #1\end{minipage}}}

\setlength{\textfloatsep}{8pt plus 2pt minus 2pt}
\setlength{\floatsep}{7pt plus 2pt minus 2pt}
\setlength{\intextsep}{7pt plus 2pt minus 2pt}
\setlength{\abovedisplayskip}{6pt plus 2pt minus 2pt}
\setlength{\belowdisplayskip}{6pt plus 2pt minus 2pt}
\setlength{\abovedisplayshortskip}{4pt plus 2pt minus 2pt}
\setlength{\belowdisplayshortskip}{4pt plus 2pt minus 2pt}
\setlist[itemize]{leftmargin=*,itemsep=1pt,topsep=2pt}


\begin{document}

\input{texfiles/title}

\input{texfiles/intro}

\input{texfiles/related}

\input{texfiles/pre}

\input{texfiles/methodology}

\input{texfiles/experiments}

\section{Conclusion}
We presented DHC, a dynamic heterogeneous context-learning framework for temporal link prediction. 
The model combines text-aware node initialization, relation-time messages, node and relation memory, recent-neighbor temporal encoding, and sequential event-memory fusion. 
On the current four datasets, DHC is competitive across scholarly citation, patent, and software temporal heterogeneous graphs. The ablations and case studies suggest that relation memory and recent-neighbor encoding should be treated as task-adaptive components. 
They are most useful when relations are sparse, non-stationary, or semantically heterogeneous.

\section*{Acknowledgment}
This work was supported by the National Natural Science Foundation of China (62272023).
\begingroup
\footnotesize
\bibliographystyle{splncs04}
\bibliography{iconip25_dhc}
\endgroup

\end{document}
